% ---
% RESUMO
% ---

\setlength{\absparsep}{18pt} % ajusta o espaçamento dos parágrafos do resumo

\begin{resumo}
% -----------------------------------------------------------------
% Estrutura recomendada (NBR 6028): um único parágrafo, terceira
% pessoa do singular, verbos na voz ativa, entre 150 e 500
% palavras. Cobrir, em ordem:
%   1. Contexto e problema.
%   2. Objetivo do trabalho.
%   3. Método adotado.
%   4. Resultados principais.
%   5. Conclusão e contribuição.
% -----------------------------------------------------------------

O rompimento de fio em fábricas de fios esmaltados de cobre e alumínio é um sintoma operacional cuja origem pode estar tanto em defeitos embarcados na matéria-prima quanto em condições degradadas do equipamento, com fenótipos de fratura frequentemente compatíveis com mais de uma hipótese e sem rotulagem padronizada de causa nos registros operacionais. Esta dissertação apresenta a qualificação do WireSense, ferramenta híbrida de apoio à decisão proposta para inferir a provável origem de cada evento de rompimento a partir da combinação de duas vertentes independentes de evidência. O método foi desenvolvido sobre dados de uma fábrica que opera com setenta e sete máquinas esmaltadeiras em paralelo, e estrutura-se em dois pilares. O Pilar de Sinais Operacionais aplica detecção de anomalias por \emph{Isolation Forest} sobre séries temporais sensoriais, com modelo treinado por máquina sobre o histórico de intervalos sem incidente. O Pilar de Histórico Estatístico opera sobre a tabela de status cruzada com as ordens de produção, e produz dois indicadores comparativos por lote: a taxa relativa ao lote, padronizada pela hora-máquina acumulada por material, e o \emph{score} de cruzamento de falhas por máquina, que mede o espalhamento dos rompimentos entre as máquinas efetivamente expostas ao lote. A integração dos dois pilares é realizada por uma regra determinística representada por duas tabelas de decisão sucessivas que produzem, primeiro, a classe de evidência histórica de matéria-prima, e em seguida o diagnóstico final entre as classes de suspeita de matéria-prima, suspeita de equipamento, investigação paralela e insuficiência de evidência. A arquitetura privilegia interpretabilidade e auditabilidade sobre otimização \emph{end-to-end}, em alternativa a modelos de fusão aprendidos que não seriam treináveis no cenário sem rótulo padronizado. Conduziu-se revisão sistemática da literatura de Computação aplicada a sistemas industriais sobre as bases IEEE Xplore, ACM Digital Library e Scopus, com seleção final de dezesseis trabalhos correlatos, que evidenciam a lacuna preenchida pela proposta. A validação retrospectiva sobre eventos com causa documentada por inspeção pericial e a implantação em regime de inferência continuada, integrada ao sistema de chão de fábrica, compõem o trabalho previsto entre esta qualificação e a defesa.

\textbf{Palavras-chave}: rompimento de fio. detecção de anomalias. estatística histórica. fusão de evidências. apoio à decisão industrial.
\end{resumo}
